\documentclass[12pt]{article}
\usepackage[utf8]{inputenc}
\usepackage[T1]{fontenc}
\usepackage{hyperref}
\title{Sourashtra Language Resources for K-12 Curriculum Development}
\author{Compiled by Claude Sonnet 5.5 (claude-sonnet-5-5) for Andr\'as Kornai}
\date{}
\begin{document}
\maketitle

\section{Introduction}
Sourashtra (Saurashtra; ISO 639-3 \texttt{saz}; Glottocode saur1248) is an Indo-Aryan language spoken by the Sourashtra community, mainly in Tamil Nadu (Madurai and surrounding towns). Sources describe it as not taught in government schools. It has its own script, which is in decline. Tamil script and Latin script are used more commonly in practice.

\section{Primary Educational Resources}
\begin{itemize}
\item Wikipedia says the script is learnt through voluntary organisations such as Sourashtra Vidya Peetam, Madurai: \url{https://en.wikipedia.org/wiki/Saurashtra_script}
\item A Deccan Chronicle item recalls that the language was taught after class hours at the Sourashtra High School in Madurai in the early 1950s. This is historical and not evidence of a current curriculum: \url{https://www.deccanchronicle.com/nation/in-other-news/090616/custodians-safeguard-saurashtra-language.html}
\item A Unicode comment document refers to a ``Saurashtra Primer'' published by Litho-press. I did not see the primer itself: \url{https://www.unicode.org/L2/L2003/03147-saurashtra-comments.pdf}
\end{itemize}

\section{Linguistic and Reference Materials}
\begin{itemize}
\item ``Semantically Classified Vocabulary of Saurashtri'' (Central Institute of Indian Languages, Mysuru), described as an aid to learning Saurashtri as a mother tongue or second language, with headwords in Tamil script: \url{https://exoticindiaart.com/book/details/semantically-classified-vocabulary-of-saurashtri-hac485}
\item Unicode proposal for the Saurashtra script, with script history and design notes: \url{https://www.unicode.org/L2/L2006/06122-saurashtra.pdf}
\item Endangered Alphabets profile of the script: \url{https://www.endangeredalphabets.net/?p=384}
\end{itemize}

\section{Religious and Community Resources}
Sourashtra Vidya Peetam, Madurai, is the community organisation named in search results for teaching the script (see above). The Endangered Alphabets profile says such groups also teach through Facebook. I found no Sourashtra Bible or dictionary in search results.

\section{Digital and Audio Resources}
The script has a Unicode block (U+A880--U+A8DF). Fonts exist, including the Suresh font used in print, as described on the Wikipedia script page. Further Unicode documents: \url{https://www.unicode.org/L2/L2005/05222r2-n2969r2-saurashtra.pdf}. I found no audio or video resources.

\section{Curriculum Development Context}
Literacy is carried by community groups, not the state. Learners often write the language in Tamil script with extra marks, so a curriculum must choose between the Saurashtra script and Tamil script. Digital support through Unicode and fonts now exists. The main gaps are graded readers, textbooks and teacher training. Collaboration with Sourashtra Vidya Peetam would be a natural starting point.
\end{document}